\documentclass[a4paper, 11pt]{article}

\usepackage[hmargin=2cm,vmargin=2.5cm,bindingoffset=0.5cm]{geometry}
\usepackage[colorlinks=true]{hyperref}
\usepackage{graphicx}
\usepackage{subfiles}
\usepackage{roboto}
\usepackage{bookmark}
\usepackage{fontspec}
\usepackage{minted}
\usepackage{color}
\usepackage{fancyhdr}
\usepackage{longtable}

% We used cascadia code before, so cc.
\setmainfont{Roboto Serif}
\setmonofont{IBM Plex Mono}
\newfontfamily\ccbold{IBM Plex Mono Bold}
\newfontfamily\ccitalic{IBM Plex Mono Italic}
\newfontfamily\ccmono{IBM Plex Mono}

\usepackage{cellspace}
\setlength\cellspacetoplimit{10pt} % Minimum space above cell content
\setlength\cellspacebottomlimit{10pt} % Minimum space below cell content

\definecolor{lightlightgray}{rgb}{0.95, 0.95, 0.95}

% italic captions
\usepackage[format=plain,
            labelfont=it,
            textfont=it]{caption}

% nicer footnotes
\usepackage[splitrule]{footmisc}

% font spacing
\setlength{\spaceskip}{0.55em plus 0.15em minus 0.15em}

% image paths
\graphicspath{{./images/}}

% paragraph spacing
\setlength{\parskip}{1em}

% no indenting paragraphs (its 2025)
\setlength{\parindent}{0pt}

% nicer paragraph headings
\usepackage{titlesec}
\titleformat{\paragraph}
  {\normalfont\bfseries}
  {}
  {0pt}
  {} 
\titleformat{\subparagraph}
  {\normalfont\bfseries}
  {}
  {0pt}
  {}

% headers
\pagestyle{fancy}
\setlength{\headheight}{12pt}
\fancyhf{}
\fancyhead[L]{\small\textbf{IGCSE Comp Sci: Arrays \& Iteration Workbook}}
\fancyhead[R]{\small Eason Qin}

\fancyfoot[C]{\thepage}

%minted stuff

\setminted{
    bgcolor=lightlightgray,
    linenos=true,
    tabsize=4
}

% small lines
\newcommand{\smalllines}{
    \begin{longtable}{p{0.9\textwidth}}
        \rule{0pt}{20pt} \\
        \hline \\
        \rule{0pt}{10pt} \\
        \hline \\
        \rule{0pt}{10pt} 
    \end{longtable}
}

% medium lines
\newcommand{\mediumlines}{
    \begin{longtable}{p{0.9\textwidth}}
        \rule{0pt}{20pt} \\
        \hline \\
        \rule{0pt}{10pt} \\
        \hline \\
        \rule{0pt}{10pt} \\
        \hline \\
        \rule{0pt}{10pt} \\
        \hline \\
        \rule{0pt}{10pt} \\
        \hline \\
        \rule{0pt}{10pt} 
    \end{longtable}
}

% large lines
\newcommand{\largelines}{
    \begin{longtable}{p{0.9\textwidth}}
        \rule{0pt}{20pt} \\
        \hline \\
        \rule{0pt}{10pt} \\
        \hline \\
        \rule{0pt}{10pt} \\
        \hline \\
        \rule{0pt}{10pt} \\
        \hline \\
        \rule{0pt}{10pt} \\
        \hline \\
        \rule{0pt}{10pt} \\
        \hline \\
        \rule{0pt}{10pt} \\
        \hline \\
        \rule{0pt}{10pt} \\
        \hline \\
        \rule{0pt}{10pt} \\
        \hline \\
        \rule{0pt}{10pt} \\
        \hline \\
        \rule{0pt}{10pt} \\
        \hline \\
        \rule{0pt}{10pt} 
    \end{longtable}
}

% larger lines
\newcommand{\verylargelines}{
    \begin{longtable}{p{0.9\textwidth}}
        \rule{0pt}{20pt} \\
        \hline \\
        \rule{0pt}{10pt} \\
        \hline \\
        \rule{0pt}{10pt} \\
        \hline \\
        \rule{0pt}{10pt} \\
        \hline \\
        \rule{0pt}{10pt} \\
        \hline \\
        \rule{0pt}{10pt} \\
        \hline \\
        \rule{0pt}{10pt} \\
        \hline \\
        \rule{0pt}{10pt} \\
        \hline \\
        \rule{0pt}{10pt} \\
        \hline \\
        \rule{0pt}{10pt} \\
        \hline \\
        \rule{0pt}{10pt} \\
        \hline \\
        \rule{0pt}{10pt} \\
        \hline \\
        \rule{0pt}{10pt} \\
        \hline \\
        \rule{0pt}{10pt} \\
        \hline \\
        \rule{0pt}{10pt} \\
        \hline \\
        \rule{0pt}{10pt} 
    \end{longtable}
}

\title{IGCSE Computer Science: Arrays \& Iteration Workbook}
\author{Eason Qin (eason@ezntek.com)}
\date{August 2026\\\emph{Second Revision}}

\begin{document}

\setlength{\arrayrulewidth}{0.3pt}

\maketitle

\section{Introduction}

{\Large \textbf{VERY IMPORTANT! This workbook only covers Iteration and Arrays (1 and 2D).}}

This document is designed to help you learn and practice Iteration (all the different types of loops) and Arrays (Array programming with both 1D and 2D arrays) in your IGCSE Computer Science syllabus. Please use the \href{https://ezntek.com/revision/pseudocode_reference.html}{Python to Pseudocode reference for help with pseudocode syntax.}

You can also use the \href{https://beancode-web.net}{Beancode Web IDE}, which ships an IGCSE Pseudocode interpreter called \href{https://github.com/ezntek/beancode}{Beancode}, which can run your pseudocode.

\textbf{NOTE:} everything in angle brackets ({\ccmono < >}) like ({\ccmono <counter>}) is a \emph{placeholder} for something; the thing you are to replace the placeholder with is described on the inside. For that example, instead of writing {\ccmono <counter>}, you write {\ccmono MyCounter}, or something similar. This document is also licensed under the CC BY-SA-NC 4.0 Int'l License.

\subsection{Prerequisites}

You need to know:

\begin{itemize}
    \item comments ({\ccmono // this is a comment})
    \item variables, constants, and pseudocode types
    \item if statements (conditionals)
    \item types of comparisons ({\ccmono >= <= > < = <>} etc.)
\end{itemize}

\newpage 
\section{Iteration}
\subfile{iteration.tex}

\newpage
\section{Arrays}
\subfile{arrays.tex}

\section{The End?}

Congratulations on coming thus far.

This is not the end (you will never stop learning new things in computer science because new things are always being made!), in terms of IGCSE, you still have to know Functions, Procedures, some builtin functions (referred to as \emph{Library Routines}), and some algorithms (sometimes referred to as \emph{standard methods of solution}) like sorting, searching, totalling, finding the minimum/maximum, etc.

Attached at the end should be an answer key; although if you have the printed copy, you may not have it.

\newpage
\section{Answers}
\subfile{answers.tex}

\end{document}
